\documentclass[a4paper,12pt]{article}

\usepackage[T2A]{fontenc}
\usepackage[utf8]{inputenc}
\usepackage[russian]{babel}
\usepackage{amsmath,amssymb,amsthm,mathtools}
\usepackage{helvet}
\renewcommand{\familydefault}{\sfdefault}
\usepackage{icomma}
\usepackage{geometry}
\usepackage{microtype}
\usepackage{tikz}
\usepackage{tkz-euclide}
\usepackage{pgfplots}
\pgfplotsset{compat=1.18}
\usepackage{tabularx,booktabs,longtable}
\usepackage{enumitem}
\usepackage{hyperref}
\usepackage{parskip}
\usepackage{fancyhdr}
\usepackage{setspace}
\usepackage{xcolor}

\geometry{left=18mm,right=18mm,top=18mm,bottom=20mm}
\onehalfspacing
\setlength{\parindent}{0pt}
\setlength{\parskip}{6pt}
\renewcommand{\arraystretch}{1.4}

\definecolor{accent}{HTML}{008080}
\definecolor{darkaccent}{HTML}{004D4D}
\definecolor{textgray}{HTML}{333333}
\definecolor{softgray}{HTML}{6E6E6E}
\definecolor{warm}{HTML}{A65D4E}
\definecolor{paper}{HTML}{F7F4EE}

\hypersetup{
  colorlinks=true,
  linkcolor=darkaccent,
  urlcolor=accent,
  pdftitle={Растворы, смеси, сплавы и проценты},
  pdfauthor={Лёвин Артём Александрович}
}

\pagestyle{fancy}
\fancyhf{}
\lhead{\textcolor{textgray}{Растворы, смеси, сплавы и проценты}}
\rhead{\textcolor{textgray}{ОГЭ по математике}}
\cfoot{\textcolor{softgray}{\thepage}}
\setlength{\headheight}{15pt}

\newcommand{\sectionline}{%
  \par\vspace{2pt}
  {\color{accent}\rule{\linewidth}{0.7pt}}
  \vspace{4pt}
}
\newcommand{\key}[1]{\textcolor{darkaccent}{\textbf{#1}}}
\newcommand{\warning}[1]{\textcolor{warm}{\textbf{#1}}}
\newcommand{\formula}[1]{%
  \begin{center}
    \begin{minipage}{0.9\linewidth}
      \centering\large\color{darkaccent}#1
    \end{minipage}
  \end{center}
}

\begin{document}

\begin{titlepage}
  \centering
  \vspace*{24mm}
  {\color{accent}\Large ЧЕК-ЛИСТ\par}
  \vspace{8mm}
  {\Huge\bfseries Растворы, смеси, сплавы\par}
  \vspace{3mm}
  {\Huge\bfseries и процентные изменения\par}
  \vspace{10mm}
  {\Large Текстовые задачи ОГЭ по математике\par}

  \vfill
  {\large Лёвин Артём Александрович\par}
  \vspace{2mm}
  {\large эксклюзивно для Софии Кальней\par}
  \vspace{8mm}
  {\large \today\par}

  \vfill
  \begin{tikzpicture}[x=1cm,y=1cm]
    \draw[accent,line width=1.4pt]
      (0,0) .. controls (2.1,1.0) and (4.2,-1.0) .. (6.3,0)
            .. controls (8.3,1.0) and (10.3,-0.8) .. (12.5,0.25);
    \fill[accent] (0,0) circle (1.6pt);
    \fill[accent] (12.5,0.25) circle (1.6pt);
  \end{tikzpicture}
  \vspace*{8mm}
\end{titlepage}

\tableofcontents
\newpage

\section{Главная идея}
\sectionline

В задачах о растворах, смесях, сплавах и высушивании удобно следить за
\key{массой выбранного компонента}. Для раствора таким компонентом обычно
является растворённое вещество; для сплава --- указанный металл; для
винограда и изюма --- сухое вещество.

\formula{
  $m_{\text{смеси}}=m_1+m_2+\dots$
  \qquad
  $m_{\text{компонента}}=m_{\text{смеси}}\cdot\omega$
}

Здесь $\omega$ --- массовая доля компонента в десятичной записи. Например,
$15\%=0{,}15$.

\subsection{Что можно складывать}

При смешивании складываются:
\begin{itemize}[leftmargin=7mm]
  \item массы растворов или смесей;
  \item массы выбранного компонента.
\end{itemize}

\warning{Концентрации напрямую не складываются.}
Итоговая концентрация находится из отношения общей массы компонента к общей
массе смеси.

\section{Массовая доля}
\sectionline

\subsection{Определение}

Массовая доля показывает, какую часть массы всей смеси составляет выбранный
компонент:

\formula{
  $\displaystyle
  \omega=\frac{m_{\text{компонента}}}{m_{\text{смеси}}}$
  \qquad или \qquad
  $\displaystyle
  \omega_{\%}=\frac{m_{\text{компонента}}}{m_{\text{смеси}}}\cdot100\%$.
}

Из основной формулы:

\formula{
  $m_{\text{компонента}}=m_{\text{смеси}}\cdot\omega$
  \qquad
  $m_{\text{смеси}}=\dfrac{m_{\text{компонента}}}{\omega}$.
}

\subsection{Простейший пример}

В $100$ г раствора содержится $10$ г соли и $90$ г воды. Тогда
\[
  \omega=\frac{10}{100}=0{,}1=10\%.
\]

В любой однородной пробе этого раствора массовая доля соли остаётся равной
$10\%$. Поэтому в $10$ г раствора содержится
\[
  10\cdot0{,}1=1\text{ г соли}.
\]

\warning{Условие однородности важно.} Если смесь разделилась на слои или
компонент выпал в осадок, такое рассуждение без дополнительных условий
неприменимо.

\section{Универсальная таблица}
\sectionline

Для большинства задач используйте три столбца:

\begin{table}[h]
\centering
\caption{Табличная модель смеси}
\begin{tabularx}{\linewidth}{@{}lXXX@{}}
\toprule
Объект & Масса смеси & Массовая доля & Масса компонента \\
\midrule
Первая смесь  & $m_1$ & $\omega_1$ & $m_1\omega_1$ \\
Вторая смесь & $m_2$ & $\omega_2$ & $m_2\omega_2$ \\
Итог         & $m_1+m_2$ & $\omega$ & $(m_1+m_2)\omega$ \\
\bottomrule
\end{tabularx}
\end{table}

Закон сохранения массы выбранного компонента даёт уравнение
\[
  m_1\omega_1+m_2\omega_2=(m_1+m_2)\omega.
\]

Отсюда:
\formula{
  $\displaystyle
  \omega=\frac{m_1\omega_1+m_2\omega_2}{m_1+m_2}$.
}

Это взвешенное среднее: раствор большей массы сильнее влияет на итоговую
концентрацию.

\section{Алгоритм решения}
\sectionline

\begin{enumerate}[leftmargin=8mm,label=\textbf{\arabic*.}]
  \item Определите, какую величину считать массой всей смеси.
  \item Выберите компонент, массу которого будете сохранять.
  \item Переведите проценты в десятичные доли.
  \item Заполните таблицу: масса смеси, доля, масса компонента.
  \item Сложите массы смесей и массы компонента.
  \item Составьте уравнение или систему уравнений.
  \item Решите её и вернитесь к вопросу задачи.
  \item Если ответ получен десятичной долей, умножьте его на $100\%$.
  \item Проверьте здравый смысл: итоговая доля при обычном смешивании должна
        лежать между исходными долями.
\end{enumerate}

\section{Разобранные модели}
\sectionline

\subsection{Смешивание двух растворов}

Смешали $4$ кг $15\%$-го раствора и $6$ кг $25\%$-го раствора.
Найдём итоговую концентрацию:
\[
  4\cdot0{,}15+6\cdot0{,}25=10x.
\]
\[
  0{,}6+1{,}5=10x,
  \qquad x=0{,}21=21\%.
\]

Проверка: $21\%$ находится между $15\%$ и $25\%$ и ближе к $25\%$,
поскольку второго раствора больше.

\subsection{Добавление чистой воды}

К $8$ кг $15\%$-го раствора добавили $4$ кг воды. Вода содержит
$0\%$ выбранного вещества:
\[
  8\cdot0{,}15+4\cdot0=12x.
\]
\[
  x=0{,}1=10\%.
\]

Масса вещества сохранилась, масса раствора увеличилась, поэтому концентрация
уменьшилась.

\subsection{Равные массы растворов}

Смешали одинаковые массы $13\%$-го и $23\%$-го растворов:
\[
  0{,}13m+0{,}23m=2mx.
\]
После сокращения на $m>0$:
\[
  x=0{,}18=18\%.
\]

Для равных масс итоговая концентрация равна обычному среднему:
\[
  \frac{13\%+23\%}{2}=18\%.
\]

\subsection{Сплавы и система уравнений}

Первый сплав содержит $20\%$ никеля, второй --- $5\%$. Из них получили
$100$ кг сплава с содержанием никеля $11\%$. Пусть массы сплавов равны
$x$ и $y$ кг:
\[
  \begin{cases}
    x+y=100,\\
    0{,}20x+0{,}05y=11.
  \end{cases}
\]

Умножим второе уравнение на $100$:
\[
  \begin{cases}
    x+y=100,\\
    20x+5y=1100.
  \end{cases}
\]

Получаем $x=40$, $y=60$. Масса первого сплава меньше массы второго на
\[
  60-40=20\text{ кг}.
\]

\subsection{Высушивание: сохраняем сухое вещество}

Виноград содержит $90\%$ влаги, значит, сухого вещества в нём $10\%$.
Изюм содержит $5\%$ влаги, значит, сухого вещества $95\%$.

Пусть для получения $20$ кг изюма требуется $x$ кг винограда. Масса сухого
вещества сохраняется:
\[
  0{,}10x=20\cdot0{,}95.
\]
\[
  x=190\text{ кг}.
\]

\key{Главный выбор:} следить нужно за тем компонентом, который не исчезает и
не добавляется в процессе.

\section{Процентное изменение}
\sectionline

\subsection{На сколько процентов изменилась величина}

\formula{
  $\displaystyle
  p=\frac{|a_{\text{нов}}-a_{\text{стар}}|}{a_{\text{стар}}}\cdot100\%$.
}

В знаменателе всегда стоит величина, \key{относительно которой} измеряют
изменение.

Например, число $20$ больше числа $12$ на
\[
  \frac{20-12}{12}\cdot100\%
  =\frac{2}{3}\cdot100\%
  =66\frac{2}{3}\%.
\]

Фразы «на сколько процентов $A$ больше $B$» и «на сколько процентов $B$
меньше $A$» дают разные ответы, потому что меняется знаменатель.

\subsection{Последовательные изменения}

Рост на $p\%$ соответствует умножению на
\[
  1+\frac{p}{100},
\]
а снижение на $p\%$ --- умножению на
\[
  1-\frac{p}{100}.
\]

Если величина $40\,000$ выросла сначала на $8\%$, затем ещё на $9\%$
от результата первого роста, то
\[
  40\,000\cdot1{,}08\cdot1{,}09=47\,088.
\]

\warning{Нельзя просто складывать проценты}, если каждое следующее изменение
считается от нового значения.

\section{Типичные ошибки}
\sectionline

\begin{enumerate}[leftmargin=8mm]
  \item Складывать концентрации вместо масс компонента.
  \item Подставлять $15$ вместо $0{,}15$.
  \item При добавлении воды приписывать ей ненулевую долю вещества.
  \item В задаче на высушивание сохранять массу воды, хотя вода испаряется.
  \item Выбирать неверный компонент для закона сохранения.
  \item Терять уравнение $x+y=m$ в задаче о двух неизвестных массах.
  \item Делить процентное изменение на новую величину вместо исходной.
  \item Складывать проценты последовательных изменений.
  \item Забывать перевести итоговую десятичную долю в проценты.
\end{enumerate}

\section{Итоговый чек-лист}
\sectionline

Перед записью ответа проверьте:
\begin{itemize}[leftmargin=8mm,label=$\square$]
  \item Я определила массу всей смеси.
  \item Я выбрала компонент, масса которого сохраняется.
  \item Я перевела проценты в десятичные доли.
  \item Я заполнила все три столбца таблицы.
  \item Я складывала массы, а не концентрации.
  \item Я учла, что доля вещества в чистой воде равна нулю.
  \item В задаче на высушивание я выбрала сухое вещество.
  \item При двух неизвестных массах я составила систему.
  \item В процентном изменении знаменатель соответствует исходной базе.
  \item Ответ соответствует смыслу и находится в допустимом диапазоне.
\end{itemize}

\newpage
\section*{Тренировочный блок}
\addcontentsline{toc}{section}{Тренировочный блок}

\subsection*{Средний уровень}
\begin{enumerate}[leftmargin=8mm]
  \item В $300$ г $12\%$-го раствора содержится некоторое вещество.
        Найдите массу этого вещества.
  \item Смешали $5$ кг $10\%$-го раствора и $3$ кг $30\%$-го раствора.
        Найдите концентрацию смеси.
  \item К $6$ кг $20\%$-го раствора добавили $2$ кг воды.
        Найдите новую концентрацию.
\end{enumerate}

\subsection*{Уровень выше среднего}
\begin{enumerate}[leftmargin=8mm,start=4]
  \item Смешали равные массы $18\%$-го и $32\%$-го растворов.
        Найдите итоговую концентрацию.
  \item К $10$ кг $8\%$-го раствора добавили чистое вещество.
        Сколько килограммов вещества нужно добавить, чтобы получить
        $20\%$-й раствор?
  \item Первый сплав содержит $30\%$ меди, второй --- $10\%$.
        Получили $80$ кг сплава с содержанием меди $18\%$.
        Найдите массы исходных сплавов и их разность.
  \item Свежие ягоды содержат $84\%$ воды, а сушёные --- $20\%$ воды.
        Сколько килограммов свежих ягод нужно для получения $12$ кг сушёных?
  \item Цена товара снизилась с $2400$ рублей до $1980$ рублей.
        На сколько процентов снизилась цена?
  \item Численность выросла с $50\,000$ сначала на $6\%$, затем
        уменьшилась на $4\%$ от нового значения. Найдите итоговую численность
        и общее процентное изменение относительно начала.
\end{enumerate}

\subsection*{Сложный уровень}
\begin{enumerate}[leftmargin=8mm,start=10]
  \item В сосуде было $12$ кг $25\%$-го раствора. Из него выпарили часть
        воды, затем добавили $2$ кг чистого вещества. Получили $40\%$-й
        раствор. Сколько килограммов воды выпарили?
\end{enumerate}

\newpage
\section*{Ответы}
\addcontentsline{toc}{section}{Ответы}

\begin{table}[h]
\centering
\caption{Ответы к тренировочному блоку}
\begin{tabularx}{\linewidth}{@{}cX@{}}
\toprule
№ & Ответ \\
\midrule
1 & $36$ г \\
2 & $17{,}5\%$ \\
3 & $15\%$ \\
4 & $25\%$ \\
5 & $1{,}5$ кг \\
6 & $32$ кг и $48$ кг; разность $16$ кг \\
7 & $60$ кг \\
8 & $17{,}5\%$ \\
9 & $50\,880$; рост на $1{,}76\%$ \\
10 & $1{,}5$ кг воды \\
\bottomrule
\end{tabularx}
\end{table}

\vfill
\begin{center}
  \textcolor{softgray}{Код самопроверки: 3615}
\end{center}

\end{document}
